% =============================================================================
% 2 FUNDAMENTAÇÃO TEÓRICA — segunda versão (18/11/2026), encorpada.
% Define os conceitos que as seções 4 a 6 usam, sem histórico exaustivo
% (skill estruturacao-introducao-logica: foco, não erudição). Acréscimos
% desta versão: agentes, integridade referencial e grão (base das classes
% de erro A e B), formas de representar o esquema, detalhes das camadas
% semânticas e o número mínimo de pares discordantes para significância.
% Citações conferidas na fonte primária em 18/11/2026 (ver referencias.bib).
% Fontes internas: docs/architecture.md (atenção, "garantia"),
% docs/semantic_layer_mapping.md (cinco categorias), 04-metodologia
% (sec:avaliacao, sec:analise). Exemplo de Q031: SQL da rodada 07a25f55
% (artifacts/raw/). Fórmulas de McNemar e Holm conferidas contra
% artifacts/reports/results_final_*-tiers.md (0:12 -> 0,00049; x3 ->
% 0,0015; 1:9 -> 0,021; x3 -> 0,064). Mínimos de pares discordantes:
% 2(1/2)^5 = 0,0625; 2(1/2)^6 = 0,031; 2(1/2)^7 = 0,0156 (x3 = 0,047).
% =============================================================================
\chapter{Fundamentação teórica}
\label{cap:fundamentacao}

Esta seção define os conceitos usados no restante do trabalho: os
modelos de linguagem e o modo como usam o contexto, os agentes construídos
sobre eles, a tarefa de tradução de linguagem natural para SQL e sua
avaliação, as formas de representar o esquema no \textit{prompt}, as
noções de integridade referencial e de grão, a camada semântica e os
testes estatísticos pareados usados na análise.

% -----------------------------------------------------------------------------
\section{Modelos de linguagem de grande escala}
\label{sec:llm}

Um modelo de linguagem de grande escala é uma rede neural treinada para
prever o próximo \textit{token} de um texto a partir dos anteriores
\cite{zhao2023survey}. O treinamento tem, em geral, duas etapas: um
pré-treinamento sobre grandes volumes de texto, em que o modelo aprende a
continuar textos, e uma adaptação posterior, em que é ajustado para seguir
instruções e responder de forma útil \cite{zhao2023survey}. É essa segunda
etapa que permite usar o modelo para uma tarefa nova apenas descrevendo a
tarefa no \textit{prompt}, sem novo treinamento. Quando o \textit{prompt}
traz só a instrução, fala-se em uso \textit{zero-shot}; quando traz também
exemplos resolvidos, em uso \textit{few-shot}, ou aprendizado em contexto
\cite{zhao2023survey}. Este trabalho usa o modelo em modo
\textit{zero-shot}.

Os modelos atuais usam a arquitetura \textit{Transformer}, cujo
componente central é o mecanismo de atenção: a cada \textit{token}
gerado, o modelo calcula pesos sobre todos os \textit{tokens} da entrada e
combina suas representações de acordo com esses pesos
\cite{vaswani2017attention}. A resposta é produzida um \textit{token} por
vez, e cada \textit{token} gerado passa a fazer parte da entrada do passo
seguinte. Nos modelos densos, todos os parâmetros participam do cálculo de
cada \textit{token}. Nos modelos de mistura de especialistas, cada camada
tem vários blocos alternativos, e apenas alguns são ativados para cada
\textit{token}, o que permite um número total de parâmetros grande com um
custo de cálculo menor. O \arq{gpt-oss-120b} é um modelo desse tipo, com
128 especialistas por camada \cite{openai2025gptoss}; o
\arq{qwen3.8-27b} é denso \cite{qwen2026qwen38}.

Tudo o que o modelo recebe em uma chamada, instruções, descrição do banco e
pergunta, forma uma única sequência de \textit{tokens}, limitada pela
janela de contexto do modelo. As interfaces de programação desses modelos
organizam a entrada em mensagens com papéis, como uma mensagem de sistema,
com instruções gerais, e mensagens de usuário, com o pedido
\cite{groq_openai}; para o modelo, porém, todas são concatenadas na mesma
sequência. Não há, nessa sequência, distinção estrutural entre uma regra e
um dado: uma instrução presente no contexto influencia a resposta na
medida em que recebe peso da atenção, e não por ser executada. Essa
propriedade é a base da distinção, usada na \autoref{sec:garantia}, entre
instrução e garantia.

A temperatura controla o grau de aleatoriedade na escolha de cada
\textit{token}: as pontuações do modelo são divididas pela temperatura
antes de virarem probabilidades, e valores baixos concentram a
probabilidade nos \textit{tokens} mais prováveis \cite{holtzman2020curious}.
No limite, a escolha se aproxima de sempre tomar o \textit{token} mais
provável. O provedor usado neste trabalho converte a temperatura zero em
um valor positivo muito pequeno ($10^{-8}$) \cite{groq_openai}. Mesmo com
temperatura baixa, servidores de inferência não garantem saídas idênticas
entre chamadas: o resultado numérico de cada passo depende de como as
requisições são agrupadas no servidor, e a carga muda esse agrupamento
\cite{he2025nondeterminism}. Em serviços acessados por API, configurações
ditas determinísticas produziram respostas e acurácias diferentes entre
execuções repetidas \cite{atil2025nondeterminism}.

Modelos de pesos abertos são aqueles cujos parâmetros treinados são
publicados e podem ser executados por terceiros, em infraestrutura
própria ou de um provedor de inferência \cite{openai2025gptoss}. Neste
trabalho, os dois modelos avaliados são de pesos abertos e foram acessados
por um provedor de inferência gratuito (\autoref{sec:modelos}).

% -----------------------------------------------------------------------------
\section{Agentes baseados em modelos de linguagem}
\label{sec:agentes}

Um agente baseado em modelo de linguagem é um sistema que usa o modelo
para decidir e executar ações em um ambiente, e não apenas para produzir
texto. Uma revisão da área organiza a construção desses agentes em quatro
módulos: perfil, que define o papel do agente; memória, que guarda
informações de interações anteriores; planejamento, que decompõe a tarefa
em passos; e ação, que transforma as decisões do modelo em operações sobre
o ambiente, como chamar ferramentas ou consultar bases de dados
\cite{wang2024agents}.

O agente deste trabalho usa o mínimo dessa estrutura, de propósito. O
perfil é a mensagem de sistema fixa; não há memória entre perguntas nem
planejamento; e a única ação é executar no banco a consulta gerada, com
até duas novas tentativas quando a consulta falha, informando ao modelo a
mensagem de erro (\autoref{sec:agente}). Com isso, o contexto fornecido no
\textit{prompt} é a única fonte de variação entre as condições.

% -----------------------------------------------------------------------------
\section{Tradução de linguagem natural para SQL}
\label{sec:text2sql}

A tradução de linguagem natural para SQL (\textit{text-to-SQL}) é a tarefa
de, dada uma pergunta em linguagem natural e o esquema de um banco de dados
relacional, produzir uma consulta SQL cujo resultado responde à pergunta
\cite{hong2024survey}. As revisões da área descrevem quatro gerações de
abordagens \cite{hong2024survey,shi2024survey,liu2025survey}: sistemas
baseados em regras e modelos de consulta escritos à mão; redes neurais do
tipo sequência para sequência, treinadas em pares de pergunta e consulta;
modelos de linguagem pré-treinados, ajustados para a tarefa; e, por fim,
modelos de linguagem de grande escala orientados por \textit{prompt}. Nesta
última geração, o modelo não é treinado sobre o banco: tudo o que ele sabe
do banco é o que está no \textit{prompt}.

A dificuldade central da tarefa está em ligar os termos da pergunta aos
elementos do esquema, como tabelas, colunas e valores, etapa chamada de
\textit{schema linking} \cite{hong2024survey}. A ligação raramente é
direta: a pergunta usa o vocabulário do negócio, e um mesmo termo pode
corresponder a mais de uma coluna ou a uma combinação de tabelas.

Um sistema de \textit{text-to-SQL} pode ser avaliado de várias formas
\cite{yu2018spider,li2023bird,shi2024survey}. A correspondência exata
compara a consulta gerada com uma consulta de referência, componente a
componente, e penaliza consultas escritas de outra forma que produzem o
mesmo resultado. A acurácia de execução executa as duas consultas na
mesma base e compara os resultados. Variações dessas métricas avaliam a
consulta em vários bancos com o mesmo esquema, para reduzir acertos por
coincidência, ou ponderam o acerto pela eficiência da consulta
\cite{li2023bird,shi2024survey}. Uma revisão observa que a
correspondência exata tende a subestimar a acurácia real, e a acurácia de
execução, a superestimá-la \cite{shi2024survey}.

As duas tendências aparecem neste trabalho. Na pergunta Q031, as
consultas \arq{SELECT SUM(total_amount) FROM orders} e
\arq{SELECT SUM(orders.total_amount) FROM orders}, geradas pelo
\arq{qwen3.8-27b} nas condições de catálogo e semântica, são textos
diferentes com o mesmo resultado, e ambas estão corretas. Já na Q011, a
consulta que conta os pedidos com cliente informado e a que conta os
pedidos ligados a um cliente existente são ambas válidas e devolvem
resultados diferentes (\autoref{sec:catalogo-vs-semantica}). Este trabalho
usa a acurácia de execução, com as normalizações descritas na
\autoref{sec:avaliacao}: uma consulta é correta quando seu resultado é
igual ao da consulta de referência.

% -----------------------------------------------------------------------------
\section{Representação do esquema no \textit{prompt}}
\label{sec:representacao-esquema}

Como o modelo só conhece o banco pelo \textit{prompt}, a forma de
descrever o esquema é parte do sistema. O \autoref{qua:representacoes}
resume as formas estudadas na literatura. Elas diferem no tipo de
informação: as primeiras descrevem a estrutura física; as descrições de
colunas acrescentam significado; os exemplos e o conhecimento por
pergunta acrescentam informação específica de cada consulta.

\begin{quadro}[htb]
  \centering
  \caption{Formas de representar o banco no \textit{prompt}}
  \label{qua:representacoes}
  \footnotesize
  \begin{tabularx}{\textwidth}{|>{\raggedright\arraybackslash}p{3.4cm}|>{\raggedright\arraybackslash}X|>{\raggedright\arraybackslash}p{3.6cm}|}
    \hline
    \textbf{Forma} & \textbf{Conteúdo} & \textbf{Estudos} \\ \hline
    Definição das tabelas & texto de criação das tabelas
      (\arq{CREATE TABLE}), com colunas, tipos e chaves &
      \citeonline{rajkumar2022evaluating}; \citeonline{gao2024dail} \\ \hline
    Lista de tabelas e colunas & nomes de tabelas e colunas em texto
      corrido, sem tipos & \citeonline{gao2024dail} \\ \hline
    Linhas de exemplo & algumas linhas de cada tabela, junto com a
      definição & \citeonline{rajkumar2022evaluating} \\ \hline
    Descrições de colunas & significado de tabelas e colunas em
      linguagem natural, como em um catálogo &
      \citeonline{li2023bird}; \citeonline{wretblad2024synthetic};
      \citeonline{bouattour2026prompting} \\ \hline
    Exemplos resolvidos & pares de pergunta e consulta
      (\textit{few-shot}) & \citeonline{gao2024dail};
      \citeonline{bouattour2026prompting} \\ \hline
    Conhecimento por pergunta & dica escrita para cada pergunta, com
      fórmulas, sinônimos ou valores & \citeonline{li2023bird};
      \citeonline{chung2025longcontext} \\ \hline
    Regras da base & convenções válidas para toda a base: grão, datas,
      relacionamentos, formato da resposta & \citeonline{rumiantsau2026semantic};
      este trabalho \\ \hline
  \end{tabularx}
  \fonte{Elaborado pelo autor a partir dos estudos citados.}
\end{quadro}

Em bancos grandes, uma etapa anterior seleciona as tabelas e colunas
relevantes para a pergunta, para que o esquema caiba no \textit{prompt}
\cite{pourreza2023din,talaei2024chess}. Cada acréscimo aumenta o número de
\textit{tokens} de entrada e, com ele, o custo e o tempo de cada chamada.

Este trabalho não usa linhas de exemplo, exemplos resolvidos, conhecimento
por pergunta nem seleção de esquema, e varia apenas o texto que descreve o
banco, em três níveis (\autoref{sec:condicoes}). O primeiro corresponde à
definição das tabelas; o segundo, às descrições de colunas; o terceiro, às
regras da base. As regras são o tipo de informação menos estudado no
\autoref{qua:representacoes} e dependem de dois conceitos apresentados a
seguir: a integridade referencial e o grão.

% -----------------------------------------------------------------------------
\section{Integridade referencial e grão}
\label{sec:integridade-grao}

No modelo relacional, os dados são organizados em relações, ou tabelas,
e as ligações entre tabelas são feitas por valores: uma coluna de uma
tabela guarda a chave de uma linha de outra \cite{codd1970relational}. Essa
coluna é chamada de chave estrangeira. A integridade referencial é a
garantia de que todo valor de uma chave estrangeira corresponde a uma
linha existente na tabela referenciada; em sistemas como o PostgreSQL, ela
é assegurada quando a restrição de chave estrangeira é declarada
\cite{postgresql16}.

Bases comerciais nem sempre têm essa garantia. Dados importados de
planilhas ou de sistemas legados podem conter referências órfãs: valores
de chave estrangeira que não correspondem a nenhuma linha da tabela
referenciada. Na base deste trabalho, a restrição não foi declarada para
preservar os dados como foram autorizados (decisão D-004), e há
referências órfãs (\autoref{sec:transformacao}). Nessa situação, duas
perguntas que parecem iguais têm respostas diferentes. ``Quantos pedidos
têm um cliente informado?'' conta as linhas em que a coluna está
preenchida (\arq{WHERE customer_id IS NOT NULL}); ``quantos pedidos estão
ligados a um cliente cadastrado?'' exige a junção com a tabela de
clientes (\arq{JOIN}), que descarta as referências órfãs. Qual das duas
uma organização chama de ``pedido com cliente associado'' é uma convenção
do negócio, e não uma propriedade do esquema. Esse é o tipo de erro
reunido na classe B (\autoref{sec:taxonomia}).

O grão de uma tabela é o que cada uma de suas linhas representa: um
pedido, um item de pedido, uma meta mensal. Na modelagem dimensional,
declarar o grão é o passo central do projeto de uma tabela de fatos, e
misturar grãos diferentes na mesma tabela ou na mesma consulta é um erro
clássico \cite{kimball2013techniques}. O erro aparece em consultas que
juntam tabelas de grãos diferentes. Se o valor total do pedido está na
tabela de pedidos e a consulta a junta com a tabela de itens, cada pedido
passa a aparecer uma vez por item, e a soma do valor total fica
multiplicada pelo número de itens. A consulta executa sem erro e devolve
um número plausível. Erros desse tipo, junto com a escolha entre colunas
de nome parecido, formam a classe A (\autoref{sec:taxonomia}).

As duas noções explicam por que o esquema físico não basta. Os nomes das
colunas e as chaves dizem como as tabelas se ligam, mas não dizem se a
ligação é completa nem em que grão cada medida deve ser somada. Essa
informação precisa vir de fora do esquema.

% -----------------------------------------------------------------------------
\section{Catálogo de dados e camada semântica}
\label{sec:camada-semantica}

Um catálogo de dados registra metadados sobre as tabelas de uma
organização, entre eles a descrição de cada tabela e coluna
\cite{ehrlinger2021catalogs,dama2017dmbok}. Em bancos relacionais, essas
descrições costumam ser mantidas como comentários associados aos objetos
(\arq{COMMENT ON TABLE}, \arq{COMMENT ON COLUMN})
\cite{postgres_comment,databricks_comment}. Uma descrição de catálogo diz
o que um objeto significa, e está presa a ele: ela pode dizer que uma
coluna guarda o valor bruto do pedido, mas não em que grão esse valor deve
ser somado nem como tratar um pedido cujo cliente não existe.

Uma camada semântica é um conjunto de definições mantido entre as tabelas
físicas e quem consulta os dados, que descreve a base em termos de negócio
e determina como ela deve ser consultada \cite{jones2024semanticlayer}. As
ferramentas industriais diferem na sintaxe, mas carregam as mesmas cinco
categorias de informação (\arq{docs/semantic_layer_mapping.md}):
\begin{alineas}
  \item definições por campo: o que cada medida e dimensão significa em
        linguagem de negócio;
  \item grão: a chave de cada tabela e o nível de agregação que ela
        representa;
  \item dimensões temporais: qual coluna de data é a data de cada medida;
  \item relacionamentos: como combinar tabelas e o que fazer quando a
        integridade referencial é imperfeita;
  \item sinônimos: a ponte entre o vocabulário dos usuários e os nomes do
        esquema.
\end{alineas}

A primeira categoria se sobrepõe ao catálogo. As outras quatro não
pertencem a nenhuma coluna isolada: são regras sobre a base como um todo.
Essa é a diferença que o desenho do experimento isola, entre o nível de
catálogo e o nível semântico.

Cada ferramenta organiza essas categorias de um modo. No dbt, o
MetricFlow define modelos semânticos com entidades, que são as chaves de
junção entre modelos, dimensões, pelas quais se agrupa, e métricas
\cite{dbt_metricflow}. No Looker, a linguagem LookML organiza o modelo em
\textit{views} com dimensões e medidas, e as junções declaram o tipo de
relacionamento entre as tabelas, como muitos para um, informação que a
ferramenta usa para calcular as medidas corretamente; grupos de dimensões
de tempo são criados de uma vez a partir de uma coluna de data
\cite{looker_lookml,looker_relationship,looker_dimgroup}. No Cube, cada \textit{cube} reúne medidas,
dimensões e junções \cite{cube_datamodel}. No Databricks, as
\textit{metric views} separam medidas de dimensões, podem ser consultadas
pelo agente da plataforma e aceitam sinônimos e nomes de exibição para
melhorar a acurácia do agente \cite{databricks_metricviews}.

Nessas ferramentas, a camada semântica é consultada por meio de um
compilador: o consumidor pede uma medida e um agrupamento, e o compilador
gera o SQL a partir das definições, inclusive as junções
\cite{dbt_slarchitecture,dbt_metricflow,looker_lookml,cube_datamodel}.
Neste trabalho, a camada semântica é apenas texto anexado ao
\textit{prompt}, e o SQL é escrito pelo modelo. Por isso, o contexto
semântico inclui também regras sobre a forma da resposta, que nas
ferramentas industriais são asseguradas pelo compilador
(\autoref{sec:industriais}).

% -----------------------------------------------------------------------------
\section{Avaliação estatística pareada}
\label{sec:estatistica}

Quando as mesmas unidades são observadas em duas condições, a comparação
deve usar testes pareados, que consideram a diferença dentro de cada
unidade \cite{agresti2013categorical}. Neste trabalho, as unidades são as
perguntas, e cada uma é respondida nas três condições.

Para um desfecho binário, como acerto ou erro, os resultados de duas
condições sobre as mesmas perguntas formam a tabela de contingência da
\autoref{tab:contingencia}. O teste de McNemar compara as duas condições
usando apenas os pares discordantes, $b$ e $c$ \cite{mcnemar1947note}.
As perguntas acertadas ou erradas nas duas condições ($a$ e $d$) não
entram no cálculo, porque não informam sobre a diferença entre elas.

\begin{table}[htb]
  \centering
  \caption{Tabela de contingência de duas condições pareadas}
  \label{tab:contingencia}
  \begin{tabular}{lcc}
    \toprule
    & Condição 2: acerto & Condição 2: erro \\
    \midrule
    Condição 1: acerto & $a$ & $b$ \\
    Condição 1: erro   & $c$ & $d$ \\
    \bottomrule
  \end{tabular}
  \fonte{Elaborado pelo autor, com base em \citeonline{agresti2013categorical}.}
\end{table}

Com $n = b + c$, sob a hipótese nula de que as duas condições têm a mesma
probabilidade de acerto, $b$ segue uma distribuição binomial com $n$
ensaios e probabilidade $1/2$. A versão exata do teste usa essa
distribuição diretamente \cite{fagerland2013mcnemar}; o valor-p bilateral
é
\begin{equation}
  p = \min\left(1,\; 2 \sum_{k=0}^{\min(b,c)} \binom{n}{k}
      \left(\frac{1}{2}\right)^{n}\right).
  \label{eq:mcnemar}
\end{equation}

Uma consequência da \autoref{eq:mcnemar} é que o número de pares
discordantes limita a significância possível. No caso mais favorável, em
que todos os $n$ pares discordantes favorecem a mesma condição, o valor-p
é $2\,(1/2)^{n}$. Com 5 pares, ele vale 0,0625 e não atinge 5\%; com 6,
vale 0,031. Assim, uma diferença só pode ser significativa a 5\% se houver
ao menos seis perguntas com resultados discordantes, e, se alguma delas
favorecer a outra condição, são necessárias mais.

Para uma medida contínua pareada, como a latência, o teste de postos
sinalizados de Wilcoxon ordena os valores absolutos das diferenças dentro
de cada par e compara a soma dos postos das diferenças positivas com a das
negativas, sem supor distribuição normal \cite{wilcoxon1945individual}.

Quando uma mesma família de dados é submetida a vários testes, a
probabilidade de ao menos um falso positivo cresce com o número de testes.
O procedimento de Holm controla essa probabilidade para a família
\cite{holm1979simple}. Ordenados os $m$ valores-p brutos em ordem crescente,
$p_{(1)} \le \dots \le p_{(m)}$, o valor-p ajustado do $i$-ésimo é
\begin{equation}
  \tilde{p}_{(i)} = \max_{j \le i}\; \min\bigl(1,\; (m - j + 1)\, p_{(j)}\bigr),
  \label{eq:holm}
\end{equation}
e cada hipótese é rejeitada quando seu valor-p ajustado é menor que o
nível de significância. O procedimento é uniformemente mais poderoso que
o ajuste de Bonferroni, que multiplica todos os valores-p por $m$
\cite{shaffer1986modified}. Neste trabalho, $m = 3$ para cada modelo: os três
pares de condições (\autoref{sec:analise}). Com o ajuste, o menor valor-p
da família é multiplicado por 3, e o mínimo de pares discordantes sobe:
com 6 pares todos na mesma direção, o valor ajustado é 0,094; com 7, é
0,047.

A acurácia de cada condição é acompanhada de um intervalo de confiança
de 95\% pela aproximação normal da proporção,
$\hat{p} \pm 1{,}96\sqrt{\hat{p}(1-\hat{p})/N}$, em que $\hat{p}$ é a
proporção de acertos e $N$ o número de perguntas
\cite{agresti1998approximate}. Com amostras pequenas, esse intervalo tem
cobertura abaixo da nominal, o que reforça a leitura conjunta com os
testes pareados.
